\documentclass[11pt]{article}
\usepackage[margin=1in]{geometry}
\usepackage{amsmath,amssymb,amsthm,booktabs,microtype}
\usepackage[colorlinks=true,linkcolor=blue!45!black,urlcolor=blue!45!black]{hyperref}
\usepackage{xcolor}
\newtheorem{theorem}{Theorem}
\newtheorem{lemma}{Lemma}
\newcommand{\Fix}{\operatorname{Fix}}
\newcommand{\Tr}{\operatorname{Tr}}
\title{A Parameter-Uniform Affine H\'enon Horseshoe:\\
Exact Coding, Stability Traces, and a Trace-Class Owner}
\author{H\'enon Route-A Working Series, C127}
\date{24 August 2026}

\begin{document}
\maketitle

\begin{abstract}
We exhibit a two-parameter family of piecewise-affine H\'enon horseshoes for
which the hyperbolic geometry, full-shift coding, every-period fixed-point
weights, and a trace-class operator identity are uniform on an explicit
rectangle.  For $(\lambda,\mu)\in[3,4]\times[1/5,1/3]$, every based binary
word of length $n$ determines exactly one periodic point and the weighted
fixed-point sum is
\[
 \frac{2^n}{(\lambda^n-1)(1-\mu^n)}.
\]
We realize this sum as the $n$th trace of a natural stability-mode operator,
derive its Fredholm product, prove a uniform zero-free disk $|z|<3/2$, and
give explicit trace-norm Lipschitz constants across the whole parameter
rectangle.  The result closes a structural uniform-parameter gate; it does
not compare the determinant with an external spectral target.
\end{abstract}

\section{Progress over the prior gate}

The prior H\'enon-track evidence included individual parameter witnesses and
finite symbolic prefixes.  Such evidence does not imply persistence on a
region.  Here the advance is a single theorem, with explicit constants, on a
two-dimensional parameter rectangle.  It simultaneously controls strip
geometry, hyperbolicity, coding, all periods, a global Hilbert-space owner,
and its determinant.  No fitted target data enter the construction.

\section{The affine horseshoe}

Let $\mathcal P=[3,4]\times[1/5,1/3]$.  For
$(\lambda,\mu)\in\mathcal P$, set
\[
 R_0=[0,\lambda^{-1}]\times[0,1],\qquad
 R_1=[1-\lambda^{-1},1]\times[0,1]
\]
and define, for $e\in\{0,1\}$,
\begin{equation}\label{eq:map}
 F_{\lambda,\mu}(x,y)=
 \bigl(\lambda x-(\lambda-1)e,\;\mu y+(1-\mu)e\bigr),
 \qquad (x,y)\in R_e.
\end{equation}
Each vertical strip maps across the square into a horizontal strip.  The
domain gap is $1-2/\lambda\ge1/3$ and the image gap is
$1-2\mu\ge1/3$.  On either branch,
\[
 DF_{\lambda,\mu}=\operatorname{diag}(\lambda,\mu),
 \qquad \lambda\ge3,\quad 0<\mu\le1/3.
\]
Thus the coordinate cones are invariant with uniform expansion and
contraction.

\begin{theorem}[uniform coding and periodic closure]\label{thm:coding}
The maximal two-sided invariant set of \eqref{eq:map} is conjugate to the
full two-shift.  A based word $e_0\ldots e_{n-1}$ has exactly one periodic
point, whose initial coordinates are
\begin{align}
 x_0&=\frac{\lambda-1}{\lambda^n-1}
       \sum_{j=0}^{n-1}\lambda^{n-1-j}e_j,\label{eq:x}\\
 y_0&=\frac{1-\mu}{1-\mu^n}
       \sum_{j=0}^{n-1}\mu^{n-1-j}e_j.\label{eq:y}
\end{align}
Consequently $\#\Fix(F^n)=2^n$, and the primitive-cycle count is
\[
 p_n=\frac1n\sum_{d\mid n}\operatorname{mobius}(d)2^{n/d}.
\]
\end{theorem}

\begin{proof}
Forward iteration of the first coordinate and backward determination of the
stable coordinate give nested intervals whose diameters are bounded by
$3^{-k}$.  The strip gaps make distinct itineraries disjoint.  Solving the
two affine closure equations gives \eqref{eq:x}--\eqref{eq:y}; these points
remain in their prescribed strips by the same nested-interval construction.
There are therefore $2^n$ based fixed points.  M\"obius inversion removes
repetitions and cyclic rooting.
\end{proof}

\section{A trace-class stability owner}

At every $q\in\Fix(F^n)$,
\[
 |\det(I-DF^n(q))|=(\lambda^n-1)(1-\mu^n).
\]
We now obtain the resulting orbit sum as an actual operator trace.  Let
\[
 \mathcal H=\mathbb C^2\otimes
 \ell^2(\mathbb N_{\ge1}\times\mathbb N_0),\qquad
 J=\begin{pmatrix}1&1\\1&1\end{pmatrix},
\]
and let $D_{\lambda,\mu}$ be diagonal with entry
$\lambda^{-r}\mu^s$ at $(r,s)$.  Define
$K_{\lambda,\mu}=J\otimes D_{\lambda,\mu}$.

\begin{theorem}[all-period trace identity]\label{thm:trace}
The operator $K_{\lambda,\mu}$ is trace class, uniformly on $\mathcal P$,
and
\begin{align}
 \|K_{\lambda,\mu}\|_1
 &=\frac{2}{(\lambda-1)(1-\mu)}\le\frac32,\label{eq:norm}\\
 \Tr K_{\lambda,\mu}^n
 &=\frac{2^n}{(\lambda^n-1)(1-\mu^n)}
 =\sum_{q\in\Fix(F^n)}\frac1{|\det(I-DF^n(q))|}.
 \label{eq:trace}
\end{align}
\end{theorem}

\begin{proof}
The nonzero eigenvalue of $J$ is $2$.  Both geometric series in
\[
 2\sum_{r\ge1}\lambda^{-r}\sum_{s\ge0}\mu^s
\]
converge uniformly, proving \eqref{eq:norm}.  Taking $n$th powers gives
$\Tr J^n=2^n$ and
\[
 \sum_{r\ge1}\lambda^{-nr}\sum_{s\ge0}\mu^{ns}
 =\frac1{(\lambda^n-1)(1-\mu^n)},
\]
which is \eqref{eq:trace}.
\end{proof}

\section{Uniform determinant control}

Trace class makes the following product an ordinary Fredholm determinant:
\begin{equation}\label{eq:fredholm}
 \det(I-zK_{\lambda,\mu})
 =\prod_{r\ge1,\,s\ge0}
   \bigl(1-2z\lambda^{-r}\mu^s\bigr).
\end{equation}
Its smallest-modulus zero is $z=\lambda/2$, so the determinant is zero-free
for $|z|<\lambda/2$ and uniformly zero-free for $|z|<3/2$.  The word
``open'' matters: at $\lambda=3$, the boundary point $z=3/2$ is a zero.

Differentiating the absolutely convergent diagonal series gives
\[
 \|\partial_\lambda K\|_1\le\frac34,
 \qquad
 \|\partial_\mu K\|_1\le\frac94
\]
on $\mathcal P$.  Integration along a coordinate path yields
\begin{equation}\label{eq:lipschitz}
 \|K_{\lambda,\mu}-K_{\lambda',\mu'}\|_1
 \le\frac34|\lambda-\lambda'|+
      \frac94|\mu-\mu'|.
\end{equation}
Equations \eqref{eq:fredholm}--\eqref{eq:lipschitz} are global parameter
statements, not neighboring-point numerical tests.

\section{Exact certificate and hostile controls}

The release evidence uses rational arithmetic on a $3\times3$ parameter grid,
checks periods through $12$, and replays representative words exactly.  An
independent implementation reconstructs every invariant.  A separate SymPy
program verifies the geometric-series and matrix identities, while sixteen
semantic mutations test gaps, counts, traces, displayed periodic coordinates,
scope flags, and route labels.

Three failure modes delimit the theorem.  If $\lambda<2$, the domain strips
can overlap; if $\mu>1/2$, the image strips can overlap.  Replacing the
stability denominator by $(\lambda\mu)^n$ is not a harmless convention: it
breaks the trace identity.  Finally, $2^n$ counts rooted words, not primitive
cycles; the latter require M\"obius inversion.

\section{Route-A assessment and limitations}

The structural uniform-parameter subgate is certified.  The strict canonical
tuple is nevertheless
\[
 (\mathrm{A1\_WEAK},\mathrm{A2\_FAIL},
  \mathrm{A3\_FAIL},\mathrm{A4\_FAIL}).
\]
A2 and A3 remain failed because there is no frozen target-facing protocol,
missing/extra-zero census, or target divisor comparison.  The internal zeros
of \eqref{eq:fredholm} must not be rebranded as arithmetic zeros.  No natural
quantization is constructed, and Route B is not authorized.  The active scope
is \texttt{NO\_BAD\_EULER\_OR\_ROOT\_NUMBER}.

\section{Conclusion}

This family supplies an explicit answer to a narrow but important Route-A
question: exact orbit coding and a trace-class stability determinant can be
made uniform on a genuine H\'enon-style parameter region.  The next separate
task is target-facing validation under a future authorized protocol; it is
not inferred from the present structural theorem.

\end{document}
